\subsection{Programación, reprogramación y cancelación de audiencias}\label{sec:impl-audiencias}

El módulo \texttt{hearings} agrega valores de dominio, consultas y un servicio de aplicación detrás de un puerto de persistencia. Cada audiencia usa un UUID estable y una revisión positiva de 32 bits. Cada escritura tiene además un UUID de operación nuevo. Programar crea la revisión uno; reemplazar o cancelar exige la revisión anterior exacta y produce su sucesor. El contador agotado se rechaza sin desbordamiento. Reutilizar una operación causa conflicto y no devuelve información del recurso que ya la empleó.

\subsubsection{Valores temporales, participantes y soporte}

La fecha conserva segundos completos y el desfase original, expresado en minutos entre $-14$ y $+14$ horas. Tanto la fecha local como su conversión UTC deben pertenecer a los años 1 a 9999. Se admiten pasado y futuro; dos representaciones del mismo instante con desfases distintos conservan valores declarados diferentes. La API exige RFC~3339 con desfase explícito, rechaza fracciones, segundos intercalares y \texttt{-00:00}, y no usa el reloj para declarar una audiencia celebrada.

Sede o conexión admiten hasta 500 caracteres Unicode; nota y motivo, hasta 1000. Se recorta espacio exterior y se convierte CRLF a LF, sin normalizar el contenido Unicode interior. No se aceptan controles ajenos al salto de línea permitido en notas y motivos. El servidor trata la conexión como texto, sin descargar su contenido. El tipo permanece fijo y los estados representan únicamente programación o cancelación.

Las referencias al directorio fijan UUID y revisión, se ordenan por UUID y rechazan duplicados aun cuando indiquen revisiones distintas. El detalle captura perfil manual o tipificado, rol, estado histórico y huellas de ficha e identidad, cuando corresponde. Una ficha editada o archivada después no reemplaza esa evidencia. En una reprogramación, las referencias retenidas deben coincidir exactamente con las de la base; las añadidas deben corresponder a cabezas activas del mismo expediente.

La individualización exige un texto de antecedente de condena y una referencia documental con UUID, versión y SHA-256. Los otros tipos rechazan ese campo. Preparar una nueva programación o reemplazo verifica el contenido cifrado, su evidencia y la admisión PDF/DOCX bajo la política de soportes, con máximo de 16~MiB. Cancelar copia el soporte histórico de la base sin volver a presentarlo como una nueva validación criptográfica o de formato. La declaración del operador no constituye una sentencia ni prueba por sí sola la existencia jurídica de una condena.

\subsubsection{Preparación, canon y recibo de operación}

La preparación devuelve comando normalizado, actor, valores, revisión resultante y dos huellas, sin reservar identificadores ni crear filas. \texttt{HEAR1} representa los valores: tipo, instante Unix y desfase, modalidad, sede, nota, participantes ordenados y antecedente opcional. Usa enteros de tamaño fijo en orden big endian, etiquetas de variantes y textos UTF-8 precedidos por su longitud en bytes. Su tamaño máximo es de 10\,726 bytes, incluyendo 32 participantes y textos de cuatro bytes por carácter. SHA-256 se calcula detrás del puerto de hash de la aplicación.

\texttt{HTXN1} vincula UUID de operación, actor, expediente y audiencia, acción, revisión esperada, expectativas de contexto, huella de valores y motivo. La cancelación no aporta expectativas administrativas o procesales nuevas: deriva su contexto de la base, mientras registra la revisión administrativa vigente por separado. El recibo no incorpora correo del autor ni hora de captura al canon y no es una firma digital. Su función es identificar un envío concreto y comprobar que la revisión consultada corresponde a aquel envío.

Confirmar repite la preparación y compara el digest esperado. La aplicación vuelve a autenticar la misma cuenta inmediatamente antes de invocar la confirmación del almacén. La persistencia revalida permisos, estado del expediente, base y fuentes después de obtener el bloqueo auditado; el instante de registro se captura después de esa espera. Alta y reemplazo exigen que el contexto esperado continúe vigente. Cancelación permite un cambio posterior de etapa, pero conserva el contexto y los participantes de su base exacta. La sesión Redis y la transacción PostgreSQL siguen siendo pasos distintos.

La familia de migraciones \texttt{0011} define raíces, revisiones, valores canónicos, recibos y restricciones de secuencia. La lectura debe reconstruir y contrastar el canon, las huellas y las fuentes históricas; no basta con devolver un JSON almacenado. Un registro inicial de etapa no tiene una huella procesal separada: el campo correspondiente es nulo, sin inventar un hash retrospectivo. Las revisiones de adopción o avance conservan su huella original. Las pruebas de integración y recuperación se distinguen de las de dominio y aplicación en la \cref{sec:pruebas-audiencias}.

\subsubsection{Contrato HTTP, consulta y autorización}

{\raggedright Las rutas del expediente parten de \rutarepo{/api/v1/cases/<case>/hearings}. Exponen contexto, preparación, alta, detalle actual, reemplazo, cancelación, revisión exacta e historia. La agenda transversal usa \rutarepo{/api/v1/hearings}. El JSON se limita a 64~KiB; se rechazan campos desconocidos o duplicados, bytes posteriores al objeto y cabeceras de tipo de contenido duplicadas. La composición comparte los mismos límites de trabajo HTTP con documentos, directorio y etapas.\par}

Owner consulta y gestiona globalmente; Litigante asignado consulta y gestiona; Paralegal asignado consulta; Cliente no accede. Un expediente cerrado mantiene las consultas autorizadas y su historia, pero rechaza mutaciones. La API distingue conflictos de revisión, contexto, participantes, operación y huella de envío, sin convertir los conflictos en reintentos automáticos. Las inconsistencias persistidas se responden como errores internos sin material sensible.

El listado por expediente ordena UUID y la historia desciende por revisión, con páginas de 1 a 100 elementos. La agenda exige un intervalo UTC positivo, de máximo 366 días y extremo final exclusivo; ordena por instante e identificador. Primero selecciona la cabeza actual, aplica autorización y filtros, y después pagina. Una revisión antigua no reaparece porque satisfaga el filtro. Los títulos y estados de los expedientes son actuales; fecha y estado de audiencia corresponden a su cabecera seleccionada.

Si se pierde una respuesta de escritura, la conciliación debe consultar la revisión exacta y comprobar audiencia, expediente, actor, operación, acción y digest esperado. La coincidencia de textos no prueba éxito. Una revisión todavía ausente conserva el resultado incierto; una revisión de otro envío exige mostrar conflicto y conservar el borrador. La preparación no congela permisos ni reserva la agenda, y no se genera otra operación para ocultar incertidumbre.

\subsubsection{Interacción de programación y agenda en Qadra}

El expediente ofrece listado, detalle, historia y editor de audiencia. El formulario distingue fecha, hora con segundos y desfase declarado; no convierte la cita al huso del navegador. Si el operador introduce únicamente horas y minutos, comunica que los segundos se registrarán como cero. La preparación presenta los valores normalizados antes de confirmar y se invalida al editar. Cambiar de expediente o sesión retira el estado correspondiente e invalida respuestas tardías.

La selección de participantes distingue las referencias históricas retenidas de las fichas vigentes que pueden añadirse. El detalle muestra los valores exactos vinculados y conserva los nombres anteriores después de una edición o archivo del directorio. Las referencias desplegables exponen UUID y revisión para distinguir fichas homónimas, sin identificar personas únicamente por el texto del nombre. La individualización añade captura explícita del antecedente y selección de versión documental; estar en Juicio no completa esos campos. Los controles de gestión consideran rol, cierre y trabajo en curso, mientras el servidor mantiene la comprobación definitiva.

La agenda es una vista transversal independiente del expediente abierto. Consulta el recurso autorizado con rango UTC, filtros y cursor conjunto de instante e identificador, sin recorrer todos los expedientes en el navegador. Agrupa las filas por día usando el desfase de visualización elegido y conserva visible la fecha, hora y desfase originales de cada cita. Abrir una fila transmite expediente, audiencia y revisión exactos; esa intención se consume una sola vez después de cargar contexto y listado, evitando iniciar esas tres consultas simultáneamente bajo el presupuesto HTTP compartido. Un cambio concurrente no convierte silenciosamente la selección en la cabecera nueva.

Ante conflictos de revisión o contexto, el editor conserva el borrador y exige consultar la base y el contexto actuales antes de volver a preparar. Si cambian participantes o soportes, exige revisar las referencias afectadas. Ante resultado incierto, conserva el envío y coteja revisión, autor, operación, acción, expectativas y huellas de valores y envío. Una ausencia temporal no permite concluir que la escritura falló. La conciliación no repite la petición; la denegación retira datos y las respuestas obsoletas no vuelven a mostrarlos. La agenda advierte que la programación puede cambiar entre consultas y que no calcula resultados, plazos o avisos.

La programación no registra resultados, asistentes efectivos, acuerdos, penas ni reparaciones, ni calcula plazos o envía alertas. La selección de una fecha o un participante no completa esas funciones. El cierre de CU-08, RF-07 y del objetivo de gestión procesal requiere evidencia adicional de esos comportamientos y de los tipos y sesiones aún no modelados.
